\chapter{Why Spreads Exist: Three Causes and the Roll Model}
\label{ch:spreads_roll}
\chaptermark{Spreads and the Roll model}

A typical US large-cap stock has a bid--ask spread as tight as
one cent on a share worth two hundred dollars --- five basis points.
Why not zero? Why not ten cents? No regulator, exchange, or single
market maker sets the width. The spread is \textbf{endogenous}: the
equilibrium output of competition between market makers posting
quotes and the order flow crossing them, determined by three
economic forces.

This chapter has two jobs. First, it names the three forces:
\emph{order-processing costs}, \emph{inventory risk}, and
\emph{adverse selection}. That three-cause framework
\citeHasbrouck{} organises
\cref{ch:spreads_roll,ch:kyle_gm,ch:market_impact,ch:market_making}.
Second, it derives the \textbf{Roll (1984) spread estimator}
\citeRoll{} from first principles: given only a tape of trade
prices --- no quotes, no volumes --- the spread leaves a recoverable
fingerprint in the lag-1 autocovariance of price changes. Four
lines of algebra and one square root suffice.

\begin{backgroundbox}[Prerequisites]
From earlier chapters:
\begin{itemize}
  \item \Cref{ch:lob}: the limit order book, resting limit orders,
  market orders versus limit orders, makers versus takers, tick size
  and lot size. The working definitions of $\bidp$, $\askp$,
  $\spreads$, and $\midp$ as quoted quantities.
  \item \Cref{ch:participants}: the \textbf{informed} vs
  \textbf{uninformed} split, and the adverse-selection toy
  calculation of \cref{subsec:toy_adverse} in which a market maker
  who quotes a zero-spread mid loses an expected $\$1.5$ per fill
  against informed counterparties.
  \item \Cref{ch:price_value}: the formal definitions of best bid,
  best ask, spread, mid, and microprice.
\end{itemize}
From outside the book:
\begin{itemize}
  \item Random variables, expectation, variance, covariance,
  independence; in particular the identities
  $\Var(X+Y) = \Var(X) + \Var(Y) + 2\Cov(X,Y)$ and
  $\Cov(X,Y) = \E[XY] - \E[X]\E[Y]$.
  \item The notion of an \textbf{iid sequence} and of a
  \textbf{random walk}. Appendix~A has a one-page refresher if any
  of these are rusty.
\end{itemize}
Programming is optional: the Roll estimator is presented as pure
arithmetic in the worked example and as a five-line Python snippet
in the Prosperity box. Either route teaches it.
\end{backgroundbox}

\section{Three reasons the spread is positive}
\label{sec:three_causes}

A market maker who quotes a positive spread is compensated for
three distinct costs. If any one were removed, the others would
still keep the spread above zero. The decomposition is
conceptual, not mechanical, and different products weight the
three components differently. Recognising which dominates in a
given market is most of what a microstructure researcher does.

\subsection{Order-processing costs}
\label{subsec:order_processing}

The most mundane component. Running a market-making operation
costs money: exchange connectivity fees, colocation rent, clearing
and settlement charges per trade, salaries of the quants and
developers maintaining the quoting system, and per-message
exchange fees for every quote update. These costs must be
amortised across the trades the market maker fills. If the spread
is narrower than the per-trade cost floor, the market maker runs
at a loss on every fill. Call this minimum width the \textbf{cost
floor}: the reason spreads on even the most contested products do
not collapse to zero.

Two counterintuitive points. First, order-processing costs are
\emph{small} on the most liquid products --- US large-cap stocks,
major currency pairs, top futures --- because volume is enormous
and fixed costs amortise over billions of trades. Second, despite
this, they often \emph{dominate} the spread on exactly those
products. Inventory and adverse selection are both much smaller on
a product like SPY (positions unwind in milliseconds; informed
flow has little edge) so the order-processing floor is what
remains, and that is why SPY's spread is stuck at one cent. On a
thinly-traded small-cap the spread is wide, but order-processing
is the \emph{smallest} of its three components.

\subsection{Inventory risk}
\label{subsec:inventory_risk}

A market maker who fills a buy now holds a long position, whether
they wanted one or not. That position is exposed to price risk
until unwound --- seconds on a liquid product, minutes on a slower
one, hours on an illiquid one. During that interval the position
can move against them: a downtick in the efficient price on a long
inventory is a real dollar loss. A rational market maker demands
compensation for that exposure, and the compensation shows up in
the spread as an additive term on both sides of the mid.

Formally, if the market maker holds inventory $\invq$ over a window
during which the efficient price has variance $\sigma^2 t$, the
variance of P\&L is $\invq^2 \sigma^2 t$, and a
mean-variance-optimising maker's compensation scales with that
variance. The Avellaneda--Stoikov treatment in
\cref{ch:market_making} derives the optimal quoting formulas and
shows how the market maker \emph{skews} quotes asymmetrically when
inventory drifts away from zero. For this chapter the qualitative
picture suffices: the spread includes a risk premium that grows
with (i) efficient-price volatility, (ii) expected holding time,
and (iii) risk aversion.

\subsection{Adverse selection}
\label{subsec:adverse_selection}

In plain English: \emph{adverse selection is the cost of getting
picked off by a counterparty who knows something you don't}. A
market maker who quotes two-sided will be preferentially hit by
traders whose information makes the quote look mispriced, so the
average fill is systematically a loser unless the quote is wide
enough to charge for the leak. This is the largest of the three
components on mid-liquidity products and the theoretical meat of
modern microstructure. Recall
\cref{subsec:toy_adverse}: a market maker faces flow with fraction
$\alpha = 0.3$ \emph{informed} traders who know the true fair value
($V = 100$ or $V = 110$), the rest uninformed. Quoting at the
unconditional mid of $105$, expected profit per fill is $-\$1.5$:
fills against informed flow are systematically on the side of the
trade that will turn out correct, so conditional on a fill the
market maker's expected P\&L is strictly negative.

The fix is to quote a \emph{spread}. An ask above the mid and a
bid below it is a tax levied on every trade. The tax must be set
high enough that gains on noise trades offset losses on informed
ones. That equilibrium spread is what the Glosten--Milgrom (1985)
and Kyle (1985) models derive in \cref{ch:kyle_gm}. The structural
observation: informed flow forces the spread strictly above zero,
and wider informed fraction means wider spread.

Adverse selection is hardest to measure directly and easiest to
get wrong in backtests. A strategy that looks profitable on
historical mid-to-mid data may pay a live adverse-selection
premium that historical mids never showed --- historical mids are
where the market quoted, not where the strategy would have been
filled. \Cref{ch:backtesting} returns to this failure mode.

\begin{keyfactbox}[The three-causes decomposition]
The quoted bid--ask spread on any market decomposes conceptually
into three strictly positive components:
\[
  \spreads \;=\; \underbrace{c_{\text{proc}}}_{\text{order-processing cost}}
         \;+\; \underbrace{c_{\text{inv}}}_{\text{inventory-risk premium}}
         \;+\; \underbrace{c_{\text{adv}}}_{\text{adverse-selection premium}}.
\]
All three are strictly positive in equilibrium. Their relative
weights depend on the product: on the most liquid products the
\emph{order-processing} floor dominates; on mid-liquidity and
illiquid products \emph{adverse selection} dominates; inventory
risk is a second-order contributor on liquid products and a
first-order one on illiquid ones. No rational market maker quotes
a spread narrower than the sum of the three.
\end{keyfactbox}

\begin{intuitionbox}
The spread is rent the market maker collects for three services.
First, they run the plumbing (exchange connections, software)
which costs money. Second, they hold a position for you when
nobody else will, charging a premium for the risk. Third, they
insure everyone against adverse selection by better-informed
counterparties. Three costs, three premiums, one observable width.
The second half of this chapter shows that \emph{even without
seeing the quotes} you can recover the total width, because all
three premiums leave the same statistical fingerprint on the trade
tape.
\end{intuitionbox}

\section{The Roll (1984) model}
\label{sec:roll_model}

Roll's puzzle, set in the early 1980s: intraday tick data was
expensive and often impossible to obtain retrospectively.
Researchers studying historical transaction costs had access only
to the \emph{tape} of printed trades --- a time-ordered sequence
of trade prices, with no quotes, quantities, or trade-direction
labels. Can you estimate the effective bid--ask spread using only
the trade-price time series?

Roll's answer: yes. The key insight is the \textbf{bid--ask
bounce}: because trades alternate between hitting the bid and
hitting the ask, trade prices zig-zag around the efficient price
even when no new information arrives, and that zig-zag leaves a
distinctive negative correlation between consecutive price
changes. Under two assumptions --- a latent random-walk
efficient price and iid trade directions --- the trade-price
series has a negative lag-1 autocovariance whose magnitude depends
only on the spread. Invert the relationship and you have a spread
estimator. The argument is four lines long.

\subsection{Setup}
\label{subsec:roll_setup}

We model two quantities. The first is an unobservable
\textbf{efficient price} $m_t$ --- the hypothetical price at which
the product would trade in a frictionless market where all
currently available information has already been incorporated.
We assume it follows a random walk:
\begin{equation}
  m_t \;=\; m_{t-1} + u_t,
  \label{eq:roll_efficient}
\end{equation}
where $u_t$ is an \iid{} shock sequence with $\E[u_t] = 0$ and
$\Var(u_t) = \sigma_u^2$. Shocks represent new information moving
the fundamental price; mean-zero means no drift on average, \iid{}
means $m_{t-1}$ tells you nothing about $u_t$.

The second quantity is the \textbf{observable trade price} $p_t$
--- what actually prints on the tape. Each trade happens either at
the bid (if it was a customer sell hitting a resting buy quote) or
at the ask (if it was a customer buy hitting a resting sell
quote). Introduce a \textbf{trade-direction indicator} $q_t$
taking the value
\[
  q_t \;=\;
  \begin{cases}
    +1 & \text{if the $t$-th trade was a customer buy (filled at the ask),} \\
    -1 & \text{if the $t$-th trade was a customer sell (filled at the bid).}
  \end{cases}
\]
The market maker posts symmetric quotes around the efficient
price: a bid at $m_t - \spreads/2$ and an ask at $m_t + \spreads/2$,
where $\spreads$ is the (constant) spread. The observable trade
price is therefore
\begin{equation}
  p_t \;=\; m_t + q_t \cdot \frac{\spreads}{2}.
  \label{eq:roll_observed}
\end{equation}
This is the key structural equation. The trade price is the
efficient price plus a \emph{bounce} of half a spread, with the
sign of the bounce determined by which side of the book the trade
hit.

We make three further assumptions, each of which is the formal
version of a convenient fiction:
\begin{enumerate}
  \item $q_t$ is symmetric iid with $\Prob(q_t = +1) = \Prob(q_t = -1)
  = 1/2$. In particular $\E[q_t] = 0$ and $\E[q_t^2] = 1$, so
  $\Var(q_t) = 1$.
  \item $\{u_t\}$ and $\{q_t\}$ are independent of each other and
  of their own pasts --- i.e.\ the efficient-price shocks carry no
  information about the direction of the current or future
  trades, and trade directions carry no information about the
  efficient-price shocks.
  \item The spread $\spreads$ is constant across the sample
  window.
\end{enumerate}
All three are false in real markets (\cref{sec:roll_when_fails}
catalogues how). But the estimator they produce remains the
workhorse for spread estimation from end-of-day data four decades
later.

The question: \emph{given only $\{p_t\}$, can we recover
$\spreads$?}

\subsection{Worked toy example, before any derivation}
\label{subsec:roll_toy}

Before the algebra, do the experiment. Pick constant efficient
price $m_t = 100.00$ (set $\sigma_u = 0$, so price changes come
from the bounce alone), spread $\spreads = 0.10$ (so $\spreads/2 =
0.05$), and a sequence of $13$ trade directions
\[
  q_1, \ldots, q_{13}
  \;=\;
  -1, +1, -1, -1, +1, -1, -1, +1, -1, -1, +1, +1, -1
\]
($5$ customer buys, $8$ customer sells, not alternating). Apply
\cref{eq:roll_observed} to get the observable prices:
\[
  p_t \;=\; 100.00 + q_t \cdot 0.05,
\]
which gives
\begin{center}
\resizebox{\linewidth}{!}{%
\begin{tabular}{rccccccccccccc}
\toprule
$t$   & $1$    & $2$    & $3$    & $4$    & $5$    & $6$    & $7$    & $8$    & $9$    & $10$   & $11$   & $12$   & $13$   \\
\midrule
$q_t$ & $-1$   & $+1$   & $-1$   & $-1$   & $+1$   & $-1$   & $-1$   & $+1$   & $-1$   & $-1$   & $+1$   & $+1$   & $-1$   \\
$p_t$ & $99.95$ & $100.05$ & $99.95$ & $99.95$ & $100.05$ & $99.95$ & $99.95$ & $100.05$ & $99.95$ & $99.95$ & $100.05$ & $100.05$ & $99.95$ \\
\bottomrule
\end{tabular}%
}
\end{center}

A tape of thirteen trade prices, all either $99.95$ or $100.05$
--- looks like noise around a round number. The differences
$\Delta p_t = p_t - p_{t-1}$ are
\begin{center}
\resizebox{\linewidth}{!}{%
\begin{tabular}{rcccccccccccc}
\toprule
$t$          & $2$    & $3$    & $4$    & $5$    & $6$    & $7$    & $8$    & $9$    & $10$   & $11$   & $12$   & $13$   \\
\midrule
$\Delta p_t$ & $+0.10$ & $-0.10$ & $0$    & $+0.10$ & $-0.10$ & $0$    & $+0.10$ & $-0.10$ & $0$    & $+0.10$ & $0$    & $-0.10$ \\
\bottomrule
\end{tabular}%
}
\end{center}

Twelve price changes. Half are zero (two fills on the same side);
the rest are $\pm 0.10$. The mean of the twelve is
\[
  \overline{\Delta p}
  \;=\; \frac{1}{12} \bigl(0.1 - 0.1 + 0 + 0.1 - 0.1 + 0 + 0.1 - 0.1 + 0 + 0.1 + 0 - 0.1\bigr)
  \;=\; \frac{0}{12}
  \;=\; 0,
\]
matching \cref{eq:roll_efficient} with $\sigma_u = 0$: no drift.
Now look at pairs of consecutive price changes. No $(+0.10,
+0.10)$ or $(-0.10, -0.10)$ pairs appear; products $\Delta p_t
\cdot \Delta p_{t-1}$ are either $0$ (one side zero) or $-0.01$
(opposite signs). \emph{None is positive.} The sample lag-1
autocovariance will be strictly negative --- Roll's point, because
a pure random walk would predict zero.

We will compute the exact number in \cref{subsec:roll_toy_payoff}
and recover the spread. But first, the general derivation.

\subsection{Derivation}
\label{subsec:roll_derivation}

Return to the general model, with $\sigma_u$ possibly nonzero, and
compute the first two moments of $\Delta p_t$. Each step
substitutes \cref{eq:roll_observed,eq:roll_efficient} and expands.

\paragraph{Step 1: rewrite $\Delta p_t$.} From
\cref{eq:roll_observed},
\begin{align*}
  \Delta p_t
  \;&=\; p_t - p_{t-1} \\
  \;&=\; \bigl(m_t + q_t \cdot \tfrac{\spreads}{2}\bigr)
       - \bigl(m_{t-1} + q_{t-1} \cdot \tfrac{\spreads}{2}\bigr) \\
  \;&=\; \underbrace{(m_t - m_{t-1})}_{=\;u_t\ \text{by \cref{eq:roll_efficient}}}
       \;+\; \frac{\spreads}{2} \,(q_t - q_{t-1}) \\
  \;&=\; u_t \;+\; \frac{\spreads}{2} \,(q_t - q_{t-1}).
\end{align*}
Price changes decompose into an efficient-price shock plus a
\emph{bounce} term proportional to the change in trade direction.
The bounce is zero when two consecutive trades hit the same side
and $\pm\spreads$ when they hit opposite sides.

\paragraph{Step 2: compute $\Var(\Delta p_t)$.} Using the identity
$\Var(X+Y) = \Var(X) + \Var(Y) + 2\Cov(X,Y)$ and the independence
of $\{u_t\}$ and $\{q_t\}$,
\begin{align*}
  \Var(\Delta p_t)
  \;&=\; \Var(u_t) \;+\; \left(\frac{\spreads}{2}\right)^{\!2} \Var(q_t - q_{t-1})
       \;+\; 2 \cdot \frac{\spreads}{2} \cdot \underbrace{\Cov(u_t,\,q_t - q_{t-1})}_{=\;0} \\
  \;&=\; \sigma_u^2 \;+\; \frac{\spreads^2}{4} \,\Var(q_t - q_{t-1}).
\end{align*}
The cross-covariance vanishes because $u_t$ is independent of the
$q$'s by assumption. Now expand $\Var(q_t - q_{t-1})$ using the
same identity and the independence of $q_t$ and $q_{t-1}$:
\[
  \Var(q_t - q_{t-1})
  \;=\; \Var(q_t) \;+\; \Var(q_{t-1}) \;-\; 2 \,\underbrace{\Cov(q_t, q_{t-1})}_{=\,0\ \text{(iid)}}
  \;=\; 1 + 1 - 0 \;=\; 2,
\]
using $\Var(q_t) = \E[q_t^2] - (\E[q_t])^2 = 1 - 0 = 1$ from the
symmetric-$\pm 1$ assumption. Therefore
\begin{equation}
  \Var(\Delta p_t) \;=\; \sigma_u^2 \;+\; \frac{\spreads^2}{2}.
  \label{eq:roll_variance}
\end{equation}
Instructive already: the variance of observable price changes
sums an information piece ($\sigma_u^2$) and a bounce piece
($\spreads^2/2$). Two unknowns; we need a second equation.

\paragraph{Step 3: compute $\Cov(\Delta p_t, \Delta p_{t-1})$.}
Substitute Step 1 for both $\Delta p_t$ and $\Delta p_{t-1}$:
\[
  \Cov(\Delta p_t, \Delta p_{t-1})
  \;=\; \Cov\!\left(\,
          u_t + \tfrac{\spreads}{2}(q_t - q_{t-1}),\;
          u_{t-1} + \tfrac{\spreads}{2}(q_{t-1} - q_{t-2})
        \,\right).
\]
Expand using bilinearity of covariance. This gives four terms:
\begin{align*}
  \Cov(\Delta p_t, \Delta p_{t-1})
  \;=\;&\; \Cov(u_t, u_{t-1}) \\
    &+\; \tfrac{\spreads}{2}\,\Cov\!\bigl(u_t,\, q_{t-1} - q_{t-2}\bigr) \\
    &+\; \tfrac{\spreads}{2}\,\Cov\!\bigl(q_t - q_{t-1},\, u_{t-1}\bigr) \\
    &+\; \tfrac{\spreads^2}{4}\,\Cov\!\bigl(q_t - q_{t-1},\, q_{t-1} - q_{t-2}\bigr).
\end{align*}
The first term vanishes: $\Cov(u_t, u_{t-1}) = 0$ because $\{u_t\}$
is \iid. The second and third terms vanish because $u$ and $q$
are independent. Only the fourth term survives:
\begin{equation}
  \Cov(\Delta p_t, \Delta p_{t-1})
  \;=\; \frac{\spreads^2}{4}\,\Cov\!\bigl(q_t - q_{t-1},\, q_{t-1} - q_{t-2}\bigr).
  \label{eq:roll_covsurvived}
\end{equation}

\paragraph{Step 4: expand the surviving $q$-covariance.} Use
$\Cov(X,Y) = \E[XY] - \E[X]\E[Y]$. Since $\E[q_t] = 0$ for all
$t$, we have $\E[q_t - q_{t-1}] = 0$ and $\E[q_{t-1} - q_{t-2}] =
0$, so both $\E$'s vanish and we need only
\[
  \E\!\bigl[(q_t - q_{t-1})(q_{t-1} - q_{t-2})\bigr].
\]
Multiply out the brackets:
\[
  (q_t - q_{t-1})(q_{t-1} - q_{t-2})
  \;=\; q_t q_{t-1} \;-\; q_t q_{t-2} \;-\; q_{t-1}^2 \;+\; q_{t-1} q_{t-2}.
\]
Now take expectations term by term, using iid of the $q$'s:
\begin{align*}
  \E[q_t q_{t-1}]   &\;=\; \E[q_t] \cdot \E[q_{t-1}] \;=\; 0 \cdot 0 \;=\; 0, \\
  \E[q_t q_{t-2}]   &\;=\; \E[q_t] \cdot \E[q_{t-2}] \;=\; 0 \cdot 0 \;=\; 0, \\
  \E[q_{t-1}^2]     &\;=\; 1 \quad\text{(since $q_{t-1} \in \{-1,+1\}$ always)}, \\
  \E[q_{t-1} q_{t-2}] &\;=\; \E[q_{t-1}] \cdot \E[q_{t-2}] \;=\; 0 \cdot 0 \;=\; 0.
\end{align*}
Summing with the signs from the product expansion:
\[
  \E\!\bigl[(q_t - q_{t-1})(q_{t-1} - q_{t-2})\bigr]
  \;=\; 0 \;-\; 0 \;-\; 1 \;+\; 0
  \;=\; -1.
\]
Substituting back into \cref{eq:roll_covsurvived}:
\begin{equation}
  \boxed{\;\Cov(\Delta p_t, \Delta p_{t-1}) \;=\; -\,\frac{\spreads^2}{4}.\;}
  \label{eq:roll_cov}
\end{equation}

\paragraph{Step 5: invert to get the estimator.} The population
lag-1 autocovariance of $\Delta p_t$ is exactly $-\spreads^2/4$ ---
a negative number whose magnitude is one quarter of the squared
spread. Solving for $\spreads$:
\begin{equation}
  \boxed{\;\spreads \;=\; 2 \,\sqrt{-\,\Cov(\Delta p_t, \Delta p_{t-1})}.\;}
  \label{eq:roll_estimator}
\end{equation}
Replace the population covariance with its sample analogue. If the
tape gives $N$ price changes $\Delta p_1, \dots, \Delta p_N$, define
$\hat c_1 \defeq \frac{1}{N}\sum_{i=1}^{N-1}
(\Delta p_{i+1} - \overline{\Delta p})(\Delta p_i - \overline{\Delta p})$
and you have the \textbf{Roll spread estimator}
\[
  \widehat{\spreads}_{\mathrm{Roll}}
  \;\defeq\; 2 \,\sqrt{-\,\hat c_1}
  \qquad (\text{undefined when } \hat c_1 \geq 0).
\]
Everything needed is in the trade-price series: no quotes, sizes,
or direction labels required.

\begin{keyfactbox}[Roll spread estimator]
Under the Roll model: (i) $m_t$ is an efficient price following
$m_t = m_{t-1} + u_t$ with $u_t$ \iid{} mean zero; (ii) the
observable trade price is $p_t = m_t + q_t \spreads/2$ with
$q_t \in \{+1,-1\}$ symmetric \iid; (iii) $\{u_t\}$ and $\{q_t\}$
are independent and the spread $\spreads$ is constant. Then
\[
  \Cov(\Delta p_t, \Delta p_{t-1}) \;=\; -\,\frac{\spreads^2}{4},
  \qquad
  \spreads \;=\; 2 \sqrt{-\,\Cov(\Delta p_t, \Delta p_{t-1})}.
\]
The Roll estimator replaces the population covariance with its
sample analogue; it is well-defined whenever the sample lag-1
autocovariance of price changes is negative.
\end{keyfactbox}

Three remarks about this result before we feed it numbers.

\begin{remark}[The efficient-price variance cancels]
\Cref{eq:roll_variance,eq:roll_cov} together say: variance of
price changes has two contributions (information and bounce) but
the lag-1 autocovariance has one --- the bounce alone. \iid{}
information innovations contribute zero lag-1 autocorrelation;
only the bounce survives at lag one, and it depends only on the
spread. Hence the Roll estimator's robustness to efficient-price
volatility: $\sigma_u^2$ dominates $\Var(\Delta p_t)$ but does not
leak into the autocovariance.
\end{remark}

\begin{remark}[Sign]
The population autocovariance is \emph{always} negative in this
model. A positive sample autocovariance signals a model-assumption
violation; the estimator is inapplicable. \Cref{sec:roll_when_fails}
lists the usual culprits.
\end{remark}

\begin{remark}[Units]
The estimator returns $\spreads$ in the same units as $p_t$.
Tape in seashells gives spread in seashells; Roll is
scale-covariant.
\end{remark}

\subsection{Applying Roll to the toy example}
\label{subsec:roll_toy_payoff}

Return to the thirteen trades of \cref{subsec:roll_toy}. The
twelve price changes were
\[
  \Delta p_2, \ldots, \Delta p_{13}
  \;=\;
  (+0.10,\;-0.10,\;0,\;+0.10,\;-0.10,\;0,\;+0.10,\;-0.10,\;0,\;+0.10,\;0,\;-0.10).
\]
The sample mean is $\overline{\Delta p} = 0$ (already computed).
Because the mean is exactly zero, the sample lag-1 autocovariance
simplifies to
\[
  \hat c_1
  \;=\; \frac{1}{N} \sum_{i=1}^{N-1} \Delta p_{t_{i+1}} \cdot \Delta p_{t_i},
\]
where $N = 12$ is the number of price changes and the $11$
products are listed below.

Write out the eleven consecutive-pair products $\Delta p_{t+1}
\cdot \Delta p_t$ explicitly:
\begin{center}
\begin{tabular}{ccr}
\toprule
Pair                         & Values            & Product \\
\midrule
$(\Delta p_2,\,\Delta p_3)$  & $(+0.10,-0.10)$   & $-0.01$ \\
$(\Delta p_3,\,\Delta p_4)$  & $(-0.10,\,0)$     & $\phantom{-}0$    \\
$(\Delta p_4,\,\Delta p_5)$  & $(0,+0.10)$       & $\phantom{-}0$    \\
$(\Delta p_5,\,\Delta p_6)$  & $(+0.10,-0.10)$   & $-0.01$ \\
$(\Delta p_6,\,\Delta p_7)$  & $(-0.10,\,0)$     & $\phantom{-}0$    \\
$(\Delta p_7,\,\Delta p_8)$  & $(0,+0.10)$       & $\phantom{-}0$    \\
$(\Delta p_8,\,\Delta p_9)$  & $(+0.10,-0.10)$   & $-0.01$ \\
$(\Delta p_9,\,\Delta p_{10})$ & $(-0.10,\,0)$   & $\phantom{-}0$    \\
$(\Delta p_{10},\,\Delta p_{11})$ & $(0,+0.10)$  & $\phantom{-}0$    \\
$(\Delta p_{11},\,\Delta p_{12})$ & $(+0.10,\,0)$ & $\phantom{-}0$    \\
$(\Delta p_{12},\,\Delta p_{13})$ & $(0,-0.10)$   & $\phantom{-}0$    \\
\midrule
\multicolumn{2}{r}{Sum}                          & $-0.03$ \\
\bottomrule
\end{tabular}
\end{center}
Three pairs contributed $-0.01$, the other eight contributed $0$;
no positive products. Sum: $-0.03$. Divide by $N = 12$:
\[
  \hat c_1 \;=\; \frac{-0.03}{12} \;=\; -0.0025.
\]
Apply \cref{eq:roll_estimator}:
\[
  \widehat{\spreads}_{\mathrm{Roll}}
  \;=\; 2 \sqrt{-\hat c_1}
  \;=\; 2 \sqrt{0.0025}
  \;=\; 2 \cdot 0.05
  \;=\; 0.10.
\]
Recovered spread is exactly $0.10$, matching the input at
\cref{subsec:roll_toy}. Job done, tape alone.

\begin{intuitionbox}
\textbf{The spread as a statistical signature of trading.}
Shown a tape of thirteen prices cold, most observers would call it
noise. Roll's observation: the noise has \emph{structure}. Every
time a trade hits the opposite side of the book, the price travels
the full spread; same side, no move (under $\sigma_u = 0$). The
bounce is a deterministic function of $q_t$, and consecutive
bounces are negatively correlated because $q_t - q_{t-1}$ is
overwhelmingly likely to reverse sign on the next step.

The fingerprint survives a genuinely random efficient price.
\iid{} shocks $u_t$ inflate $\Var(\Delta p_t)$ but contribute
nothing at lag one. Spread-signal and information-noise live in
orthogonal statistics; Roll's estimator reads the statistic that
isolates the spread.
\end{intuitionbox}

\section{When Roll works, and when it doesn't}
\label{sec:roll_when_fails}

Roll's derivation relied on three assumptions: independent-increment
efficient price, iid symmetric trade directions, constant spread.
Each fails in real markets in a specific way, biasing the estimator
in a specific direction.

\subsection{Assumption 1: trade directions are not iid}

In real markets, trade directions cluster. A parent order worked
through a VWAP algorithm generates a long run of buys, followed by
a long run of sells when the next institutional order arrives.
Runs produce long sequences of zero $\Delta p$ punctuated by
jumps; the sample lag-1 autocovariance is pushed toward zero,
biasing the Roll estimate downward. The Hawkes-process treatment
in \cref{ch:order_flow} formalises this and gives corrections.

\subsection{Assumption 2: the efficient price is not
independent of order flow}

Informed flow moves the efficient price in the direction of the
trade: an informed buy at the ask is followed, on average, by a
permanent upward drift as beliefs update. This permanent impact
creates positive correlation between the bounce and the next
efficient-price shock, spilling into the lag-1 autocovariance with
opposite sign. The effects partly cancel, and Roll under-reports
the true spread whenever adverse selection is material. The
Glosten--Milgrom and Kyle models of \cref{ch:kyle_gm} formalise
this, and the Hasbrouck VAR decomposition corrects for it given
quote-level data.

\subsection{Assumption 3: the spread is not constant}

Spreads widen around news, open, and close; narrow midday; widen
when large makers step back. Roll applied across a regime change
returns a time-average, biased toward whichever regime dominates
the sample.

\begin{pitfallbox}[Don't apply Roll to illiquid data]
Roll fails worst on illiquid products with few prints. A stock
trading three times an hour on a day with real-information drift
has $\Delta p_t$ observations dominated by $u_t$ shocks, not the
bounce --- so $\sigma_u^2 \gg \spreads^2/2$ in
\cref{eq:roll_variance}. Population lag-1 autocovariance is still
$-\spreads^2/4$, but the sample version from a handful of
drift-dominated changes can easily come out positive by sampling
noise, leaving the estimator undefined (\texttt{NaN}).

Worse, an analyst who flips the sign by hand --- ``surely it must
be negative, take the absolute value'' --- computes a statistic
with no theoretical backing. Roll is a \emph{model}, and on
illiquid data the model is misspecified. When sample autocov comes
out positive, discard, don't patch. For illiquid products use the
Corwin--Schultz high-low estimator, or direct quote data if
available.
\end{pitfallbox}

These caveats make Roll a \emph{starting point}, not the end.
Model-light, quick, requires only trade prices. Modern practice
uses it as a sanity check and default when nothing better is
available, cross-validated against Hasbrouck or direct quotes when
those exist.

\section{A Prosperity application}
\label{sec:roll_prosperity}

The Roll estimator transfers cleanly to Prosperity, where trades
are the simplest data channel and quotes (though available through
\texttt{TradingState}) can be noisy with competitor-algo signalling.
Below: Roll on a simulated \texttt{RAINFOREST\_RESIN} tape (success)
and a drifting-\texttt{KELP} tape (failure).

\begin{prosperitybox}
\textbf{Roll on \texttt{RAINFOREST\_RESIN} (stationary, iid
directions).} Prosperity's \texttt{RAINFOREST\_RESIN} trades
around a near-constant fair value of about $10{,}000$ seashells,
quoted with a small market-maker spread of $\spreads = 0.40$
seashells $(s/2 = 0.20)$. Take a stretch of $13$ consecutive
trade prints with an efficient price of $10{,}000$ throughout
$(\sigma_u = 0)$ and the same trade-direction sequence as
\cref{subsec:roll_toy}:
\[
  q \;=\; (-1,+1,-1,-1,+1,-1,-1,+1,-1,-1,+1,+1,-1).
\]
Using $p_t = 10{,}000 + q_t \cdot 0.20$, the trade prices are
\[
  9999.8,\;10000.2,\;9999.8,\;9999.8,\;10000.2,\;9999.8,\;9999.8,
\]
\[
  10000.2,\;9999.8,\;9999.8,\;10000.2,\;10000.2,\;9999.8.
\]
The twelve price changes are
\[
  \Delta p \;=\; (+0.4,-0.4,\,0,+0.4,-0.4,\,0,+0.4,-0.4,\,0,+0.4,\,0,-0.4).
\]
The mean price change is $0$ exactly (the sequence is balanced).
Each non-zero product of consecutive price changes is $-0.16$;
three such pairs contribute, all others are zero, so
\[
  \sum_{t=2}^{N-1} \Delta p_{t+1}\,\Delta p_t
  \;=\; 3 \cdot (-0.16) \;=\; -0.48,
  \qquad
  \hat c_1 \;=\; \frac{-0.48}{12} \;=\; -0.04.
\]
The Roll estimator returns
\[
  \widehat{\spreads}_{\mathrm{Roll}}
  \;=\; 2\sqrt{0.04} \;=\; 2 \cdot 0.2 \;=\; 0.40,
\]
recovering the true spread exactly. The happy case:
\texttt{RAINFOREST\_RESIN} is stationary with a tight, near-constant
spread and balanced flow. Roll is the right tool.

\medskip
\textbf{Roll on a drifting \texttt{KELP} tape (fails).} Now suppose
\texttt{KELP} is trending upward during the sample window --- its
efficient price drifting by $+0.5$ seashells per trade --- and
nine consecutive taker orders arrive on the ask (so
$q_t = +1$ throughout). With a true spread of $\spreads = 1$, the
trade prices are
\begin{gather*}
  2000.5,\;2001.0,\;2001.5,\;2002.0,\;2002.5, \\
  2003.0,\;2003.5,\;2004.0,\;2004.5,\;2005.0,
\end{gather*}
giving nine price changes $\Delta p_t = +0.5$ for every $t$. The
sample mean of $\Delta p$ is $0.5$, and every centred change
$(\Delta p_t - 0.5)$ is exactly zero, so
\[
  \hat c_1 \;=\; \frac{1}{N}\sum_t (\Delta p_{t+1} - 0.5)(\Delta p_t - 0.5)
  \;=\; 0.
\]
Plugging into \cref{eq:roll_estimator}: $\widehat{\spreads}_{\mathrm{Roll}}
= 0$, on a product with true spread $1$. Trending plus one-sided
flow destroyed both assumptions at once; Roll returned garbage.

\medskip
\textbf{Lesson.} Roll fits stationary, two-sided, uninformed-flow
regimes like \texttt{RAINFOREST\_RESIN}; it fails on trending
one-sided regimes like \texttt{KELP} during a directional burst.
For the latter, compute the spread directly from the top of book
in \texttt{TradingState}.

\medskip
\textbf{Reference implementation.} A five-line NumPy version of
the Roll estimator, suitable for dropping into any Prosperity
research notebook:
\begin{lstlisting}[language=Python, caption={Roll spread estimator from trade prices.}]
import numpy as np

def roll_spread(prices: np.ndarray) -> float:
    """Estimate the effective spread from a trade-price series."""
    dp = np.diff(prices)                    # N price changes
    dp_c = dp - dp.mean()                   # centre once, full-sample
    c1 = (dp_c[1:] * dp_c[:-1]).sum() / len(dp)   # lag-1 autocov, /N
    if c1 >= 0:
        return float("nan")                 # Roll undefined
    return 2.0 * np.sqrt(-c1)
\end{lstlisting}
Feed it the RESIN array and it returns $0.40$; feed it KELP and
it returns \texttt{nan}, matching the hand arithmetic. Subtlety:
\texttt{np.cov} with \texttt{bias=True} would differ because it
re-centres each slice by its own mean and divides by $N-1$. The
implementation above matches the chapter's derivation exactly.
\end{prosperitybox}

\begin{gsbox}
On a sell-side desk, Roll is a workhorse for historical spread
estimation from end-of-day OHLC on products without licensed
intraday tick data. A PM asking ``what was the effective spread on
the Mexican peso non-deliverable forward in February 2019?''
typically gets a Roll-based estimate unless the desk paid for the
tick archive --- and for exotic products, usually they haven't.

On HFT desks with abundant tick data, Roll is superseded by the
Hasbrouck VAR decomposition, which separates permanent
(information) and transient (noise) impacts and recovers Kyle's
$\lamk$ plus the adverse-selection share of the spread. But for
daily/weekly backtests and execution-assumption validation, Roll
remains the default. \Cref{ch:backtesting} covers walk-forward
validation: estimate the spread on the trailing window, apply it
as cost on the next, iterating through the sample.
\end{gsbox}

\section{Key facts to memorise}
\label{sec:ch4_keyfacts}

\begin{itemize}
  \item \textbf{Three-causes decomposition.} The bid--ask spread
  decomposes into \emph{order-processing costs} (the mechanical
  floor), \emph{inventory risk} (compensation for holding a
  position), and \emph{adverse selection} (insurance against
  informed flow). All three strictly positive; their relative
  weights depend on the product's liquidity.
  \item \textbf{Roll's model.} Trade prices $p_t = m_t + q_t
  \spreads/2$, where $m_t$ is a latent random-walk efficient price
  and $q_t \in \{+1,-1\}$ is an iid trade-direction indicator.
  \item \textbf{Roll's result.} $\Cov(\Delta p_t, \Delta p_{t-1}) =
  -\spreads^2/4$, so
  $\spreads = 2\sqrt{-\Cov(\Delta p_t, \Delta p_{t-1})}$.
  Derivable from the tape alone; no quote data required.
  \item \textbf{Why it works.} The $\sigma_u^2$ information-shock
  variance dominates $\Var(\Delta p_t)$ but contributes \emph{zero}
  to the lag-1 autocovariance; the spread dominates the lag-1
  autocovariance and nothing else. The two statistics are
  orthogonal.
  \item \textbf{Assumptions required.} Random-walk efficient price,
  iid symmetric trade directions, constant spread. All three are
  wrong in real markets in specific, correctable ways.
  \item \textbf{When to use Roll.} Stationary markets, balanced
  two-sided order flow, tape data only. Default for daily-bar
  backtests on products without tick data.
  \item \textbf{When not to use Roll.} Trending markets, illiquid
  products with few prints, one-sided order flow, news-event
  windows, and \emph{any} sample in which the computed $\hat c_1$
  comes out positive (which means the model is misspecified; the
  correct response is to discard, not to patch).
  \item \textbf{Connection to \cref{ch:participants}.} The
  adverse-selection component of the three-cause decomposition is
  exactly the $\$1.5$-per-fill loss derived in the toy calculation
  of \cref{subsec:toy_adverse}: informed flow imposes a tax, and
  the spread is the market maker's mechanism for collecting that
  tax back from everyone.
  \item \textbf{Cross-reference map.} Adverse selection in detail:
  \cref{ch:kyle_gm}. Inventory-risk premium in detail:
  \cref{ch:market_making}. Clustered order flow and the first
  assumption: \cref{ch:order_flow}. Using Roll inside a
  historical backtest: \cref{ch:backtesting}.
\end{itemize}
